\section{Related Work}
\label{chap:related-work}

Since little research exists on mailbox types, we discuss related work with a focus on mechanizations of the $\pi$-calculus and session types \cite{hondaTypesDyadicInteraction1993}.
These are tailored to languages that make use of \emph{communication channels} rather than mailboxes.

Ambal et al. \cite{ambalHOPiCoq2021} present a Rocq formalization of a higher-order $\pi$-calculus, where messages are able to carry processes.
Their work focuses on the representation of binders in name restrictions of processes.
The authors provide four versions of their mechanization, each with a different representation of variables.
% locally nameless, de Bruijn indices, nominal representation (where variables are by convention distinct from each other) and higher-order abstract syntax.
They conclude that their de Bruijn formalization is the most concise, but requires complex manipulation of de Bruijn indices and intricate proofs about renaming lemmas.
This fact can also be observed in our work.
% While de Bruijn indices remove the need for $\alpha$-equivalence, their maintenance can become complex.

Goto et al. \cite{gotoExtensibleApproachSession2016} provide a mechanization of polymorphic session types for the $\pi$-calculus in Rocq.
Their mechanization uses the locally-nameless \cite{chargueraudLocallyNamelessRepresentation2012} approach for variables.
% In their mechanization, they use the locally-nameless representation for variables in the calculus.
Similar to our work, the authors first provide their definitions for processes and types independently.
% In turn, one can argue about untyped, as well as typed processes.
They are able to prove that their type system guarantees preservation and safety properties.

Thiemann~\cite{thiemannIntrinsicallyTypedMechanizedSemantics2019}
presents a mechanization of session types in Agda.
In contrast to other mechanizations, his approach makes use of
\emph{intrinsically} typed terms, meaning that
only well-typed terms can be constructed.
As is common with intrinsically typed mechanizations, variables are represented by de Bruijn indices \cite{wadlerProgrammingLanguageFoundations2022a}.
% The operational semantics are modeled by two layers.
% At the top layer, a scheduler selects processes from a thread pool.
% Then, at the bottom layer, single expressions are evaluated.
% Evaluation happens in the style of an abstract machine.
The preservation and progress properties are proven by construction.
Intrinsic typing can lead to reduced proof sizes and fewer lemmas \cite{wadlerProgrammingLanguageFoundations2022a}.
Whether such a representation would also work for Pat is worth exploring.

Tirore et al.~\cite{tiroreMultipartyAsynchronousSession2025} provide a mechanization of multiparty asynchronous session types in Rocq.
They rely on the original definitions by Honda et al.~\cite{hondaMultipartyAsynchronousSession2016}.
During their mechanization they discovered flaws with the original definitions, such that type preservation does not hold in that system.
They propose a new type system for which preservation and other properties hold, formally proving this claim.
To manage the complexity of variables, they use the \emph{Autosubst2} library \cite{starkMechanisingSyntaxBinders2019}, which generates Rocq code based on the de Bruijn representation.
In contrast to our formalization, the authors choose to represent environments as a list of pairs, with name and type.
They claim that this representation, in particular, simplified context splitting for parallel composition.

Zalakain and Dardha \cite{zalakainLeftoversMechanisationAgda2021} present a mechanization of a type system for the $\pi$-calculus in Agda.
Their approach facilitates \emph{leftover typing}, where the typing judgment includes a second leftover environment.
The leftover environment contains types not used in the typing of the term, leaving them available for parallel processes.
This technique eliminates the need for environment splits but would likely be
difficult to use due to the need for type combination.
